% Generated draft: requires mathematical and source review.
\documentclass[11pt]{article}
\usepackage[margin=25mm]{geometry}
\usepackage{fontspec}
\setmainfont{Latin Modern Roman}
\usepackage{amsmath,amssymb,mathtools,unicode-math}
\setmathfont{Latin Modern Math}
\usepackage{enumitem,needspace,xurl,hyperref}
\setlength{\parindent}{0pt}
\setlength{\parskip}{5pt}
\setlength{\emergencystretch}{3em}
\newcommand{\N}{\mathbb{N}}
\newcommand{\Z}{\mathbb{Z}}
\newcommand{\Q}{\mathbb{Q}}
\newcommand{\R}{\mathbb{R}}
\newcommand{\C}{\mathbb{C}}
\newcommand{\F}{\mathbb{F}}
\newcommand{\A}{\mathbb{A}}
\newcommand{\PP}{\mathbb{P}}
\newcommand{\E}{\mathbb{E}}
\newcommand{\eps}{\varepsilon}
\newcommand{\dd}{\,\mathrm{d}}
\begin{document}
\section*{UnsolvedMath ID: 800956}
\textbf{Source problem code:} catalogue entry
\subsection{Definitions and mathematical statement}
Let \(M\) be a closed manifold, meaning compact and without boundary, with \(\dim M\geq 3\). The supplied record denotes its space of metrics by \(\mathcal{G}(M)\), but does not specify the regularity of the metrics or the topology on this space. Under the intended Riemannian interpretation, with sufficient regularity to define a geodesic flow, set
\[
S_gM=\{(x,v)\in TM:g_x(v,v)=1\},
\]
and let \(\varphi_g^t:S_gM\to S_gM\), \(t\in\R\), be the geodesic flow. Call this flow transitive precisely when
\[
\exists z\in S_gM\qquad
\overline{\{\varphi_g^t(z):t\in\R\}}^{\,S_gM}=S_gM.
\]
Thus the source's notation \(SM\) is interpreted as the unit tangent bundle for the metric under consideration.

Write
\[
\mathcal{T}(M)=\{g\in\mathcal{G}(M):\varphi_g^t\text{ has a dense orbit in }S_gM\}.
\]
The question asks whether, for every closed \(M\) of dimension at least three,
\[
\overline{\mathcal{T}(M)}^{\,\tau}=\mathcal{G}(M),
\]
where \(\tau\) must be the intended topology on \(\mathcal{G}(M)\). Equivalently, does every nonempty \(\tau\)-open subset of \(\mathcal{G}(M)\) contain a metric with transitive geodesic flow?

The supplied excerpt leaves \(\mathcal{G}(M)\) and \(\tau\) undefined, so this is a parameterized formulation rather than a fully specified density assertion. It also does not state a connectedness hypothesis on \(M\); connectedness has not been added here. These qualifications require clarification from the source context.
\subsection{Short English statement}
For a closed manifold of dimension at least three, must metrics whose geodesic flow has a dense orbit be dense among all metrics? The metric class, its topology, and any intended connectedness hypothesis are not specified in the supplied excerpt.
\subsection{Sources}
\begin{itemize}
\item Local UnsolvedMath registry: permanent ID \(800956\), source problem code \texttt{catalogue entry}, title \emph{Open Problems and Questions about Geodesics --- Question 36}. The registry attributes the question to Burns and Matveev and records the proposed year as 2019.
\item Supplied source: Burns, Matveev, \emph{Open Problems and Questions about Geodesics} (2019), Question 36, \url{https://arxiv.org/abs/1308.5417}. The supplied provenance records access on 2026-10-16 and extraction by \texttt{arxiv\_latex}; these are supplied provenance claims, not independently verified access or source checks.
\item The registry labels the problem \texttt{open}. This is a supplied status label, not an independently verified assessment of the literature.
\end{itemize}
\end{document}
